%% ============================================================
%%  preamble.tex  -  Cấu hình chung cho toàn bộ báo cáo
%%  Biên dịch bằng XeLaTeX (latexmk -xelatex main.tex)
%%  CHỈ NGƯỜI QUẢN LÝ REPO ĐƯỢC SỬA FILE NÀY.
%%  Cần thêm gói/lệnh mới -> nhắn người quản lý.
%% ============================================================

%% ---------- 1. Thông tin đề tài (dùng lại ở trang bìa, PDF metadata) ----------
\newcommand{\tendetai}{XÂY DỰNG HỆ THỐNG CUNG CẤP NƯỚC SẠCH QUY MÔ NHỎ CHO MỘT CỤM DÂN CƯ TẠI KHU VỰC TÂN PHÚ ĐÔNG, TỈNH TIỀN GIANG}

%% ---------- 2. Font + tiếng Việt ----------
\usepackage{fontspec}
% Dùng Times New Roman nếu máy có (compile local), không có thì dùng
% TeX Gyre Termes (tương đương Times, có sẵn trên GitHub Actions).
\IfFontExistsTF{Times New Roman}{%
  \setmainfont{Times New Roman}%
}{%
  \setmainfont{TeX Gyre Termes}%
}
\setsansfont{TeX Gyre Heros}
\setmonofont{TeX Gyre Cursor}[Scale=MatchLowercase]

% polyglossia (không dùng babel) vì babel tiếng Việt không có sẵn trên mọi bản TeX Live
\usepackage{polyglossia}
\setmainlanguage{vietnamese}
% Đổi tên một số mục tự sinh cho khớp mẫu báo cáo
\addto\captionsvietnamese{%
  \renewcommand{\contentsname}{Mục lục}%
  \renewcommand{\listfigurename}{Danh mục hình ảnh}%
  \renewcommand{\listtablename}{Danh mục bảng}%
  \renewcommand{\bibname}{Tài liệu tham khảo}%
  \renewcommand{\figurename}{Hình}%
  \renewcommand{\tablename}{Bảng}%
}

%% ---------- 3. Bố cục trang ----------
\usepackage[a4paper, top=2.5cm, bottom=2.5cm, left=3cm, right=2cm]{geometry}
\usepackage{setspace}
\setstretch{1.3}                 % giãn dòng 1.3
\usepackage{indentfirst}         % thụt đầu dòng cả đoạn đầu tiên
\setlength{\parindent}{1cm}
\setlength{\parskip}{4pt plus 1pt}

%% ---------- 4. Màu, bảng biểu, danh sách ----------
\usepackage[table,dvipsnames]{xcolor}
\usepackage{array}
\usepackage{tabularx}
\usepackage{longtable}           % bảng dài nhiều trang (WBS, ngân sách...)
\usepackage{multirow}
\usepackage{makecell}
\usepackage{booktabs}
\usepackage{adjustbox}           % co bảng/hình quá rộng: \adjustbox{max width=\textwidth}{...}
\renewcommand{\arraystretch}{1.25}

\usepackage{enumitem}
\setlist{nosep, leftmargin=*}
\setlist[itemize]{topsep=3pt, itemsep=2pt}
\setlist[enumerate]{topsep=3pt, itemsep=2pt}

%% ---------- 5. Toán, đơn vị ----------
\usepackage{amsmath, amssymb}
\usepackage{siunitx}
\sisetup{
  output-decimal-marker = {,},   % 1,5 (dấu phẩy thập phân)
  group-separator       = {.},   % 1.000.000 (dấu chấm ngăn cách hàng nghìn)
  group-minimum-digits  = 4,
  per-mode              = symbol,
}
\DeclareSIUnit{\ngay}{ngày}
\DeclareSIUnit{\ho}{hộ}
\DeclareSIUnit{\nguoi}{người}
\newcommand{\VND}{\,VNĐ}
% Ví dụ: \SI{200}{\liter\per\ngay}, \SI{50}{\ho}, \num{150000000}\VND

%% ---------- 6. Hình ảnh ----------
\usepackage{graphicx}
\graphicspath{{figures/}}        % ghi: \includegraphics{part3/gantt.png}
\usepackage{float}               % dùng được [H]
\usepackage{pdflscape}           % trang ngang: \begin{landscape}...\end{landscape}
\usepackage[font=small, labelfont=bf, labelsep=period]{caption}
\usepackage{subcaption}

%% ---------- 7. Vẽ sơ đồ: WBS, Gantt, mạng công việc ----------
\usepackage{tikz}
\usetikzlibrary{arrows.meta, positioning, shapes.geometric, calc, fit, backgrounds}
\usepackage{pgfgantt}            % biểu đồ Gantt
\usepackage[edges]{forest}       % cây phân rã công việc (WBS)
\usepackage{pgffor}

%% ---------- 8. Tiêu đề chương/mục (theo mục lục: "2.1. Tên mục") ----------
\usepackage{titlesec}
\titleformat{\chapter}[block]
  {\centering\bfseries\fontsize{16}{20}\selectfont}
  {\MakeUppercase{\chaptername}~\thechapter:}{0.5em}
  {\MakeUppercase}
\titleformat{name=\chapter, numberless}[block]
  {\centering\bfseries\fontsize{16}{20}\selectfont}
  {}{0pt}{\MakeUppercase}
\titlespacing*{\chapter}{0pt}{-10pt}{24pt}

\titleformat{\section}
  {\normalfont\bfseries\fontsize{14}{17}\selectfont}{\thesection.}{0.5em}{}
\titleformat{\subsection}
  {\normalfont\bfseries\fontsize{13}{16}\selectfont}{\thesubsection.}{0.5em}{}
\titleformat{\subsubsection}
  {\normalfont\bfseries\itshape\fontsize{12}{15}\selectfont}{\thesubsubsection.}{0.5em}{}
\titlespacing*{\section}{0pt}{14pt plus 2pt}{6pt}
\titlespacing*{\subsection}{0pt}{10pt plus 2pt}{4pt}
\titlespacing*{\subsubsection}{0pt}{8pt plus 2pt}{3pt}

\setcounter{secnumdepth}{3}      % đánh số đến 3.2.1.1 nếu cần
\setcounter{tocdepth}{2}         % mục lục hiện đến 3.2.1

%% ---------- 9. Đầu trang / chân trang ----------
\usepackage{fancyhdr}
\renewcommand{\chaptermark}[1]{\markboth{\chaptername\ \thechapter.\ #1}{}}
\pagestyle{fancy}
\fancyhf{}
\fancyhead[L]{\small\itshape\leftmark}
\fancyfoot[C]{\small\thepage}
\renewcommand{\headrulewidth}{0.4pt}
\renewcommand{\footrulewidth}{0pt}
\fancypagestyle{plain}{%          % trang đầu chương: chỉ có số trang
  \fancyhf{}
  \fancyfoot[C]{\small\thepage}
  \renewcommand{\headrulewidth}{0pt}
}

%% ---------- 10. Liên kết, URL, PDF metadata ----------
\usepackage{xurl}                % cho phép ngắt dòng URL dài
\usepackage[unicode]{hyperref}
\hypersetup{
  colorlinks = true,
  linkcolor  = black,
  citecolor  = black,
  urlcolor   = blue!50!black,
  pdftitle   = {Xây dựng hệ thống cung cấp nước sạch quy mô nhỏ cho một cụm dân cư tại khu vực Tân Phú Đông, tỉnh Tiền Giang},
  pdflang    = {vi},
}
\usepackage{bookmark}

%% ---------- 11. Lệnh tiện ích dùng chung ----------
% Đánh dấu chỗ chưa viết xong (hiện chữ đỏ, dễ tìm): \todo{bổ sung số liệu}
\newcommand{\todo}[1]{\textcolor{red}{\textbf{[TODO: #1]}}}

%% ---------- 12. (TÙY CHỌN) Tài liệu tham khảo bằng biblatex ----------
%% Bỏ dấu % ở khối dưới nếu cả nhóm dùng file .bib riêng (bib/part1.bib ... part7.bib),
%% rồi thêm \printbibliography trước \end{document} trong main.tex.
%% Mỗi người đặt khóa trích dẫn có tiền tố phần, ví dụ: \cite{p3-pmbok}.
%
% \usepackage[backend=biber, style=numeric, sorting=none]{biblatex}
% \foreach \i in {1,...,7}{%
%   \IfFileExists{bib/part\i.bib}{\addbibresource{bib/part\i.bib}}{}%
% }
